\documentclass[10pt,a4paper]{amsart}
\usepackage{amsmath,amssymb,amsthm,graphicx}
\usepackage{verbatim}
\newenvironment{reference}{\comment}{\endcomment}
%\newenvironment{reference}{}{}
\newenvironment{slogan}{\comment}{\endcomment}
\newenvironment{history}{\comment}{\endcomment}

% For commutative diagrams we use Xy-pic
\usepackage[all]{xy}

% We use 2cell for 2-commutative diagrams.
\xyoption{2cell}
% Theorem environments.
%
\theoremstyle{plain}
\newtheorem{theorem}[subsection]{Theorem}
\newtheorem{proposition}[subsection]{Proposition}
\newtheorem{lemma}[subsection]{Lemma}

\theoremstyle{definition}
\newtheorem{definition}[subsection]{Definition}
\newtheorem{example}[subsection]{Example}
\newtheorem{exercise}[subsection]{Exercise}
\newtheorem{situation}[subsection]{Situation}

\theoremstyle{remark}
\newtheorem{remark}[subsection]{Remark}
\newtheorem{remarks}[subsection]{Remarks}

\numberwithin{equation}{subsection}

% Macros
%
\def\lim{\mathop{\mathrm{lim}}\nolimits}
\def\colim{\mathop{\mathrm{colim}}\nolimits}
\def\Spec{\mathop{\mathrm{Spec}}}
\def\Hom{\mathop{\mathrm{Hom}}\nolimits}
\def\Ext{\mathop{\mathrm{Ext}}\nolimits}
\def\SheafHom{\mathop{\mathcal{H}\!\mathit{om}}\nolimits}
\def\SheafExt{\mathop{\mathcal{E}\!\mathit{xt}}\nolimits}
\def\Sch{\mathit{Sch}}
\def\Mor{\mathop{\mathrm{Mor}}\nolimits}
\def\Ob{\mathop{\mathrm{Ob}}\nolimits}
\def\Sh{\mathop{\mathit{Sh}}\nolimits}
\def\NL{\mathop{N\!L}\nolimits}
\def\CH{\mathop{\mathrm{CH}}\nolimits}
\def\proetale{{pro\text{-}\acute{e}tale}}
\def\etale{{\acute{e}tale}}
\def\QCoh{\mathit{QCoh}}
\def\Ker{\mathop{\mathrm{Ker}}}
\def\Im{\mathop{\mathrm{Im}}}
\def\Coker{\mathop{\mathrm{Coker}}}
\def\Coim{\mathop{\mathrm{Coim}}}

% Boxtimes
%
\DeclareMathSymbol{\boxtimes}{\mathbin}{AMSa}{"02}

%
% Macros for moduli stacks/spaces
%
\def\QCohstack{\mathcal{QC}\!\mathit{oh}}
\def\Cohstack{\mathcal{C}\!\mathit{oh}}
\def\Spacesstack{\mathcal{S}\!\mathit{paces}}
\def\Quotfunctor{\mathrm{Quot}}
\def\Hilbfunctor{\mathrm{Hilb}}
\def\Curvesstack{\mathcal{C}\!\mathit{urves}}
\def\Polarizedstack{\mathcal{P}\!\mathit{olarized}}
\def\Complexesstack{\mathcal{C}\!\mathit{omplexes}}
% \Pic is the operator that assigns to X its picard group, usage \Pic(X)
% \Picardstack_{X/B} denotes the Picard stack of X over B
% \Picardfunctor_{X/B} denotes the Picard functor of X over B
\def\Pic{\mathop{\mathrm{Pic}}\nolimits}
\def\Picardstack{\mathcal{P}\!\mathit{ic}}
\def\Picardfunctor{\mathrm{Pic}}
\def\Deformationcategory{\mathcal{D}\!\mathit{ef}}
\usepackage{cleveref}
\crefname{originallemma}{original lemma}{original lemmas}
\Crefname{originallemma}{Original lemma}{Original lemmas}
\makeatletter
\newlabel{coherent-lemma-flat-base-change-cohomology}{{02KH}{source}{}{sourcepointer-1}{}}
\newlabel{coherent-lemma-flat-base-change-cohomology@cref}{{[originallemma][0][]02KH}{[1][0][]source}{}{}{}}
\makeatother
\crefname{originalproposition}{original proposition}{original propositions}
\Crefname{originalproposition}{Original proposition}{Original propositions}
\makeatletter
\newlabel{cohomology-proposition-vanishing-Noetherian}{{02UZ}{source}{}{sourcepointer-2}{}}
\newlabel{cohomology-proposition-vanishing-Noetherian@cref}{{[originalproposition][0][]02UZ}{[1][0][]source}{}{}{}}
\makeatother
\crefname{originallemma}{original lemma}{original lemmas}
\Crefname{originallemma}{Original lemma}{Original lemmas}
\makeatletter
\newlabel{pic-lemma-torsion-descends}{{0CD5}{source}{}{sourcepointer-3}{}}
\newlabel{pic-lemma-torsion-descends@cref}{{[originallemma][0][]0CD5}{[1][0][]source}{}{}{}}
\makeatother
\crefname{originalreference}{original reference}{original references}
\Crefname{originalreference}{Original reference}{Original references}
\makeatletter
\newlabel{proposition-torsion-picard-reduced-proper}{{0C20}{source}{}{sourcepointer-4}{}}
\newlabel{proposition-torsion-picard-reduced-proper@cref}{{[originalreference][0][]0C20}{[1][0][]source}{}{}{}}
\makeatother
\crefname{originallemma}{original lemma}{original lemmas}
\Crefname{originallemma}{Original lemma}{Original lemmas}
\makeatletter
\newlabel{varieties-lemma-change-fields-pic}{{0CC5}{source}{}{sourcepointer-5}{}}
\newlabel{varieties-lemma-change-fields-pic@cref}{{[originallemma][0][]0CC5}{[1][0][]source}{}{}{}}
\makeatother
\title{l-torsion in the Picard group of a proper curve is bounded by the first cohomology plus the geometric genus}
\author{The Stacks Project Authors}
\date{stacks-project:curves.tex:local-snapshot-20261012; GNU FDL 1.2 or later}
\begin{document}
\maketitle
\noindent\textit{Editorial note.} All delivered mathematical text is the author's own, copied from the pinned Stacks Project chapter file curves.tex; no statement or proof step was rewritten, reordered or removed.\par\medskip
\noindent\textit{Editorial note.} The counted proof is the author's complete argument: the base-change reduction, the H\textasciicircum{}1 surjection onto the reduction with its vanishing justification, the component and torsion isomorphisms, and the final proposition citation are all retained verbatim.\par\medskip
\noindent\textit{Editorial note.} The label lemma-bound-torsion-simple names the counted statement, and the proof environment immediately following it in the source is delivered verbatim as its proof.\par\medskip
\noindent\textit{Editorial note.} The counted statement is carried by the author's own lemma environment in the chapter source.\par\medskip
\noindent\textit{Allowed standard prerequisites (editorial).}\begin{itemize}
\item Picard and cohomology under base change. Varieties, Lemma varieties-lemma-change-fields-pic; Cohomology of Schemes, Lemma coherent-lemma-flat-base-change-cohomology. Injects Pic(X) into Pic(X\_\{overline\{k\}\}) and preserves dim H\textasciicircum{}1 under the field base change, reducing to the algebraically closed case.
\item Reduction to the reduced curve. Cohomology, Proposition cohomology-proposition-vanishing-Noetherian; Picard Schemes of Curves, Lemma pic-lemma-torsion-descends. The reduction surjection stays surjective on H\textasciicircum{}1 by Cech vanishing, and X\_red to X is an isomorphism on components and on l-torsion of Picard groups.
\item The reduced-proper torsion proposition. Same chapter, Proposition proposition-torsion-picard-reduced-proper. Supplies the final bound dim Pic(X)[ell] <= dim H\textasciicircum{}1(O\_X) + g\_geom for reduced proper curves over algebraically closed fields.
\end{itemize}\medskip
\ifdefined\hypertarget\hypertarget{sourcepointer-1}{}\fi
\noindent\textit{Source reference.} Original lemma 02KH: sources/stacks-chow/tags.txt, source lines 3372--3372. Prerequisite pointer invoked at proof line(s) 3972; it names a standard fact of a sibling chapter that the counted argument uses, and every step of the counted proof built on it is retained verbatim in the delivered body.\par\medskip
\ifdefined\hypertarget\hypertarget{sourcepointer-2}{}\fi
\noindent\textit{Source reference.} Original proposition 02UZ: sources/stacks-chow/tags.txt, source lines 3750--3750. Prerequisite pointer invoked at proof line(s) 3983; it names a standard fact of a sibling chapter that the counted argument uses, and every step of the counted proof built on it is retained verbatim in the delivered body.\par\medskip
\ifdefined\hypertarget\hypertarget{sourcepointer-3}{}\fi
\noindent\textit{Source reference.} Original lemma 0CD5: sources/stacks-chow/tags.txt, source lines 15530--15530. Prerequisite pointer invoked at proof line(s) 3987; it names a standard fact of a sibling chapter that the counted argument uses, and every step of the counted proof built on it is retained verbatim in the delivered body.\par\medskip
\ifdefined\hypertarget\hypertarget{sourcepointer-4}{}\fi
\noindent\textit{Source reference.} Original reference 0C20: sources/stacks-curves/pinned-original-source.tex, source lines 3701--3701. Prerequisite pointer invoked at proof line(s) 3989; it names a standard fact of the same chapter that the counted argument uses, and every step of the counted proof built on it is retained verbatim in the delivered body.\par\medskip
\ifdefined\hypertarget\hypertarget{sourcepointer-5}{}\fi
\noindent\textit{Source reference.} Original lemma 0CC5: sources/stacks-chow/tags.txt, source lines 15495--15495. Prerequisite pointer invoked at proof line(s) 3971; it names a standard fact of a sibling chapter that the counted argument uses, and every step of the counted proof built on it is retained verbatim in the delivered body.\par\medskip
\begin{lemma}
\label{lemma-bound-torsion-simple}
Let $k$ be a field. Let $X$ be a proper scheme of dimension $\leq 1$ over $k$.
Let $\ell$ be a prime number invertible in $k$. Then
$$
\dim_{\mathbf{F}_\ell} \Pic(X)[\ell] \leq
\dim_k H^1(X, \mathcal{O}_X) + g_{geom}(X/k)
$$
where $g_{geom}(X/k)$ is as defined above.
\end{lemma}

\begin{proof}
The map $\Pic(X) \to \Pic(X_{\overline{k}})$
is injective by Varieties, Lemma \ref{varieties-lemma-change-fields-pic}.
By Cohomology of Schemes, Lemma \ref{coherent-lemma-flat-base-change-cohomology}
$\dim_k H^1(X, \mathcal{O}_X)$ equals
$\dim_{\overline{k}} H^1(X_{\overline{k}}, \mathcal{O}_{X_{\overline{k}}})$.
Hence we may assume $k$ is algebraically closed.

\medskip\noindent
Let $X_{red}$ be the reduction of $X$. Then the surjection
$\mathcal{O}_X \to \mathcal{O}_{X_{red}}$ induces a surjection
$H^1(X, \mathcal{O}_X) \to H^1(X, \mathcal{O}_{X_{red}})$
because cohomology of quasi-coherent sheaves vanishes in degrees
$\geq 2$ by
Cohomology, Proposition \ref{cohomology-proposition-vanishing-Noetherian}.
Since $X_{red} \to X$ induces an isomorphism on irreducible
components over $\overline{k}$ and an isomorphism on
$\ell$-torsion in Picard groups
(Picard Schemes of Curves, Lemma \ref{pic-lemma-torsion-descends})
we may replace $X$ by $X_{red}$. In this way we reduce to
Proposition \ref{proposition-torsion-picard-reduced-proper}.
\end{proof}

\end{document}
